% =====================================================================
% Projeto de dissertação — PPGEco/UFES
%
% Estrutura: ABNT NBR 15287 (projeto de pesquisa), classe abntex2.
% Compilar:  make projeto      (ou: cd projeto && latexmk -pdf main.tex)
%
% Este documento é uma CONDENSAÇÃO de paper/ (versão longa do Ensaio 1) e dos
% documentos de decisão em docs/. Todo número vem de lá; a origem de cada um
% está no Apêndice B. A bibliografia é a mesma de paper/ (fonte única).
%
% RASCUNHO x VERSÃO FINAL. Com \rascunhotrue o PDF mostra, em cor:
%   \decisao{n}{...}  pontos que dependem do aval do orientador
%                     (pauta em docs/pauta-orientador-2026-11.md)
%   \pendente{...}    informação que falta (data de defesa, dados de capa)
% Para a versão de entrega, trocar por \rascunhofalse: as marcas somem.
% =====================================================================
\documentclass[12pt,oneside,a4paper,chapter=TITLE,english,brazil]{abntex2}

\usepackage{lmodern}
\usepackage[T1]{fontenc}
\usepackage[utf8]{inputenc}
\usepackage{indentfirst}
\usepackage{microtype}
\usepackage{amsmath,amssymb}
\usepackage{booktabs}
\usepackage{graphicx}
\usepackage{xcolor}
\usepackage{pdflscape}
\usepackage[alf]{abntex2cite}

% --- Rascunho ou versão final -----------------------------------------
\newif\ifrascunho
\rascunhotrue
% \rascunhofalse

\newcommand{\decisao}[2]{\ifrascunho{\color{red!65!black}\textbf{[Decisão D#1:} #2\textbf{]}}\fi}
\newcommand{\pendente}[1]{\ifrascunho{\color{blue!60!black}\textbf{[A informar:} #1\textbf{]}}\fi}

% --- Ponte com paper/: as seções de lá usam natbib ----------------------
% \citet{chave} -> Autor (ano);  \citep{chave} -> (AUTOR, ano)
\newcommand{\citet}[1]{\citeonline{#1}}
\newcommand{\citep}[1]{\cite{#1}}

% --- Quadros (ABNT distingue quadro de tabela) ---------------------------
\newcommand{\quadroname}{Quadro}
\newcommand{\listofquadrosname}{Lista de quadros}
\newfloat{quadro}{loq}{\quadroname}
\newlistof{listofquadros}{loq}{\listofquadrosname}
\newlistentry{quadro}{loq}{0}
\counterwithout{quadro}{chapter}
\renewcommand{\cftquadroname}{\quadroname\space}
\renewcommand*{\cftquadroaftersnum}{\hfill--\hfill}

% --- Notação -----------------------------------------------------------
\newcommand{\ATT}{ATT(d\,|\,d)}

% --- Dados de capa e folha de rosto ------------------------------------
\titulo{Ensaios sobre a proibição da pulverização aérea de agrotóxicos no
  Ceará: avaliação de impacto e escolha de instrumento}
\autor{Felipe Carvalho Souza Santos}
\local{Vitória}
\data{2026}
\orientador{Prof. Dr. Renato Seixas}
\instituicao{%
  Universidade Federal do Espírito Santo
  \par
  Centro de Ciências Jurídicas e Econômicas
  \par
  Programa de Pós-Graduação em Economia}
\tipotrabalho{Projeto de pesquisa}
\preambulo{Projeto de dissertação apresentado ao Programa de Pós-Graduação em
  Economia da Universidade Federal do Espírito Santo, como requisito parcial
  para a obtenção do título de Mestre em Economia.}

\makeatletter
\hypersetup{
  pdftitle={\@title},
  pdfauthor={\@author},
  pdfsubject={\imprimirpreambulo},
  colorlinks=true,
  linkcolor=black,
  citecolor=black,
  urlcolor=black,
  bookmarksdepth=4
}
\makeatother

% --- Capa com a instituição no topo (guia de normalização da UFES) -------
\renewcommand{\imprimircapa}{%
  \begin{capa}%
    \center
    {\large\imprimirinstituicao}
    \vfill
    {\large\imprimirautor}
    \vfill
    \begin{center}
      \ABNTEXchapterfont\bfseries\LARGE\imprimirtitulo
    \end{center}
    \vfill
    {\large\imprimirlocal}\par
    {\large\imprimirdata}
    \vspace*{1cm}
  \end{capa}}

\setlength{\parindent}{1.25cm}
\setlength{\parskip}{0.2cm}

\begin{document}
\selectlanguage{brazil}
\frenchspacing

% ----------------------------------------------------------------------
% Elementos pré-textuais
% ----------------------------------------------------------------------
\imprimircapa
\imprimirfolhaderosto

\begin{siglas}
  \item[ACR] \emph{Average causal response}, a derivada da curva dose-resposta
  \item[ADAGRI] Agência de Defesa Agropecuária do Estado do Ceará
  \item[ADI] Ação Direta de Inconstitucionalidade
  \item[ATT] \emph{Average treatment effect on the treated}
  \item[DATASUS] Departamento de Informática do Sistema Único de Saúde
  \item[DiD] Diferenças em diferenças
  \item[FAO-GAEZ] \emph{Global Agro-Ecological Zones}, da Organização das Nações
    Unidas para a Alimentação e a Agricultura
  \item[IBGE] Instituto Brasileiro de Geografia e Estatística
  \item[LAI] Lei de Acesso à Informação
  \item[MAPA] Ministério da Agricultura e Pecuária
  \item[MDE] Efeito mínimo detectável (\emph{minimum detectable effect})
  \item[PAM] Produção Agrícola Municipal
  \item[SEMACE] Superintendência Estadual do Meio Ambiente do Ceará
  \item[SIM] Sistema de Informações sobre Mortalidade
  \item[SINAN] Sistema de Informação de Agravos de Notificação
  \item[SINASC] Sistema de Informações sobre Nascidos Vivos
  \item[SISAGUA] Sistema de Informação de Vigilância da Qualidade da Água para
    Consumo Humano
  \item[SPT] Tendências paralelas fortes (\emph{strong parallel trends})
  \item[STF] Supremo Tribunal Federal
  \item[VPP] Valor preditivo positivo
\end{siglas}

\pdfbookmark[0]{\contentsname}{toc}
\tableofcontents*
\cleardoublepage

% ----------------------------------------------------------------------
% Elementos textuais
% ----------------------------------------------------------------------
\textual

% =====================================================================
% 1 Introdução — tema, problema, hipóteses, objetivos, justificativa
% Fontes: paper/secoes/01-introducao.tex; docs/justificativa.md §§1–5.
% =====================================================================

\chapter{Introdução}
\label{cap:introducao}

\section{Tema e problema de pesquisa}
\label{sec:tema}

Em 2010, José Maria Filho, o Zé Maria do Tomé, foi assassinado na Chapada do
Apodi, no interior do Ceará. Ele denunciava havia anos a pulverização aérea de
agrotóxicos sobre os perímetros irrigados da região, onde aviões sobrevoam
plantações de banana e melão vizinhas de casas, escolas e açudes de
abastecimento. A denúncia não tratava do veneno aplicado na lavoura, e sim do
veneno que não fica na lavoura. Nove anos depois, o estado do Ceará proibiu a
prática em todo o território, pela Lei Estadual nº 16.820/2019, a primeira
proibição estadual do gênero no Brasil, e deu à lei o nome dele.

A pulverização aérea é um caso de externalidade negativa em que o mecanismo
físico é visível. Parte da calda liberada pela aeronave não atinge o alvo: é
carregada pelo vento e se deposita fora da área tratada, fenômeno que a
literatura técnica chama de deriva. Pela conta que a própria Assembleia
Legislativa usou para decidir, reproduzida na justificativa do Projeto de Lei
nº 18/2015 e ali atribuída à Embrapa, 19\% do produto se dispersa para fora da
área de aplicação. Quem contrata a aplicação paga pelo serviço e pelo produto.
Quem mora a favor do vento arca com o restante.

A lei cearense responde a esse problema com o instrumento mais forte
disponível, a proibição do método. O Supremo Tribunal Federal a julgou
constitucional em 2023 (ADI 6137), a Lei Estadual nº 19.135/2024 abriu exceção
para aeronaves remotamente pilotadas, e nova ação segue pendente (ADI 7794). A
política está em vigor, em disputa e sujeita a ser copiada por outros estados,
e as decisões sobre mantê-la, revogá-la ou substituí-la vêm sendo tomadas sem
estimativa de efeito.

Esta dissertação se organiza em dois ensaios, e cada um responde a uma
pergunta.

\begin{description}
  \item[Ensaio 1.] Qual é o efeito causal da proibição da pulverização aérea de
    agrotóxicos no Ceará sobre desfechos de saúde ao nascer, e como esse efeito
    varia com a intensidade com que cada município estava exposto à prática
    antes da proibição?
  \item[Ensaio 2.] Dado o objetivo que a lei declara, a proibição é o
    instrumento adequado em comparação com a taxação e com a regulação por
    zonas-tampão, e que informação seria necessária para ordenar esses
    instrumentos?
\end{description}

\section{Hipóteses}
\label{sec:hipoteses}

As hipóteses do Ensaio 1 foram fixadas por escrito antes de qualquer
estimação (Apêndice~\ref{ap:preespecificacao}). Elas não apontam todas na mesma
direção, e é isso que torna o conjunto informativo.

\begin{description}
  \item[H1 (principal).] Se a pulverização aérea gera externalidade de saúde por
    deriva, a proibição reduz a exposição da população e melhora os desfechos
    perinatais, e a melhora é maior onde a exposição anterior era maior.
  \item[H2 (seleção para nascimento vivo).] Se a proibição reduz o óbito fetal,
    fetos que antes não chegavam a nascer passam a compor a amostra pela cauda
    inferior de peso, e o efeito sobre o peso médio aparece atenuado ou
    invertido. Por isso o óbito fetal é o segundo desfecho confirmatório.
  \item[H3 (substituição, de sinal oposto).] Proibido o avião, a aplicação migra
    para trator e pulverizador costal. A exposição da população cai, e a
    intoxicação aguda ocupacional sobe com a intensidade anterior de aplicação.
  \item[H4 (falsificação).] A lei alcança o método, não a molécula, e a
    disponibilidade do produto não muda. A intoxicação autoprovocada, que
    responde à disponibilidade, não deve se mover com a exposição anterior.
\end{description}

O Ensaio 2 não testa hipótese empírica própria. Ele parte de uma proposição a
examinar: a de que a comparação entre proibir e taxar pelo critério de
bem-estar esperado \cite{weitzman1974} exige objetos que o Ensaio 1 não
identifica, e de que o critério coerente com o objetivo declarado da lei é o de
custo-efetividade e confiabilidade na obtenção de exposição aérea nula
(Capítulo~\ref{cap:ensaio2}).

\section{Objetivos}
\label{sec:objetivos}

\subsection{Objetivo geral}

Avaliar a proibição estadual da pulverização aérea de agrotóxicos no Ceará como
resposta a uma externalidade negativa, estimando seu efeito sobre desfechos de
saúde ao nascer e examinando a escolha do instrumento regulatório.

\subsection{Objetivos específicos}

\begin{alineas}
  \item Construir um painel município-mês para o Ceará que reúna nascimentos
    (SINASC), óbitos fetais (SIM), notificações de intoxicação (SINAN) e
    intensidade agrícola anterior à proibição (PAM/IBGE);
  \item estimar o efeito da proibição sobre o peso ao nascer, o baixo peso e a
    prematuridade por diferenças em diferenças com tratamento contínuo
    \cite{callaway2024continuous}, com inferência adequada a poucos grupos;
  \item medir o poder estatístico do desenho antes e depois da estimação e
    reportar o resultado contra critério de falsificação fixado previamente;
  \item medir a distância entre a variável de exposição disponível, que é área
    plantada, e a prática que a lei proibiu, que é a aplicação por aeronave, e
    obter registros administrativos que reduzam essa distância;
  \item avaliar desenhos alternativos que concentrem a comparação onde havia
    pulverização aérea a remover, com destaque para o contraste sazonal entre
    meses de pico e de entressafra de aplicação;
  \item delimitar, no Ensaio 2, quais objetos a escolha entre proibição, taxa e
    zona-tampão exigiria, quais deles a evidência disponível informa e sob que
    critério os instrumentos podem ser comparados.
\end{alineas}

\section{Justificativa}
\label{sec:justificativa}

A literatura causal sobre agrotóxicos e saúde ao nascer mediu três coisas, e
nenhuma delas é uma proibição de método.

Mediu substâncias pelo canal hídrico. \citet{dias2023down} exploram a
expansão da soja transgênica no Brasil e a direção do fluxo d'água dentro da
bacia, e documentam piora dos desfechos ao nascer nas populações a jusante.
Mediu substâncias pela difusão de uma tecnologia. \citet{reynier2025glyphosate}
aplicam lógica semelhante nos Estados Unidos, com aptidão agroclimática como
fonte de variação, e estimam redução de 23 a 32 g no peso ao nascer à
intensidade média de exposição ao glifosato, com efeito doze vezes maior no
decil inferior de peso esperado do que no superior. E mediu um método aéreo em
circunstância distinta: \citet{camacho2017health} estudam a aspersão aérea de
herbicida sobre cultivos ilícitos na Colômbia, onde a pulverização foi
introduzida como política de Estado, ao passo que no Ceará ela foi retirada.

Três lacunas justificam o trabalho.

A primeira é de objeto. O tratamento aqui é um ato regulatório com data, texto
legal e reversão parcial, que proíbe o modo de aplicar sem nomear princípio
ativo algum. As buscas feitas no corpus acadêmico e nos periódicos brasileiros
de economia não retornaram avaliação causal de proibição de método de
aplicação. A afirmação cabível é que a busca não encontrou, e não que o estudo
não exista.

A segunda é de contexto químico. Os desenhos de referência são construídos
sobre glifosato em culturas transgênicas. Na Chapada do Apodi, a tabela de
análises de água reproduzida na justificativa do projeto de lei registra
procimidona e carbaril em 23 de 23 pontos de coleta, carbofurano em 18,
fenitrotiona em 16 e glifosato em 4. Trata-se de fungicida e inseticida sobre
fruticultura irrigada. O análogo mais próximo é \citet{calzada2023bananas}, que
estudam a fumigação aérea de bananais no Equador e encontram déficit de peso ao
nascer sob exposição alta.

A terceira é de política. \citet{greenstone2014india} mostram, para a Índia,
que regulação ambiental no papel não equivale a redução de exposição. Saber se
a proibição cearense alterou desfechos de saúde é pergunta em aberto, com
consequência direta para os estados que discutem medida semelhante e para o
Ceará, que em 2024 trocou a proibição total por regulação de desempenho para
drones.

Há ainda uma justificativa de método, que os resultados preliminares tornaram
central. O conjunto de desfechos que os registros administrativos brasileiros
oferecem é amplo, e o poder estatístico deste desenho é curto. Essa combinação
produz achado espúrio com aparência de rigor. O projeto fixou desfecho,
exposição e critério de falsificação antes de estimar, e reporta o que o
desenho distingue e o que não distingue. Um resultado inconclusivo com a medida
de quanto falta para concluir é informação útil sobre uma política em disputa.

\section{Estado do projeto e estrutura do documento}
\label{sec:estado}

Este projeto não parte do zero. O pipeline de dados e de estimação está
implementado e versionado, as bases principais foram adquiridas, o plano de
análise foi fechado em 21 de setembro de 2026 e a estimação do Ensaio 1 foi
executada uma vez com dados reais. O Capítulo~\ref{cap:resultados} reporta
esses resultados como preliminares. Eles mostram que o desenho original mede o
efeito de forma reduzida da proibição segundo uma \emph{proxy} de bananicultura
e que a precisão obtida é inferior ao critério fixado. A etapa seguinte, que
este projeto propõe, é melhorar a medida do tratamento e testar um desenho com
mais poder.

O Capítulo~\ref{cap:contexto} apresenta a lei e o contexto.
O Capítulo~\ref{cap:referencial} traz o referencial teórico e a literatura
empírica. O Capítulo~\ref{cap:metodologia} descreve a metodologia do Ensaio 1.
O Capítulo~\ref{cap:resultados} reporta os resultados preliminares.
O Capítulo~\ref{cap:ensaio2} expõe a proposta do Ensaio 2.
O Capítulo~\ref{cap:plano} apresenta o plano de trabalho e os riscos, e o
Capítulo~\ref{cap:cronograma}, o cronograma.

% =====================================================================
% 2 Contexto institucional e legal
% Fontes: paper/secoes/02-background.tex; docs/legislacao/README.md.
% =====================================================================

\chapter{Contexto institucional e legal}
\label{cap:contexto}

\section{A Lei Zé Maria do Tomé}
\label{sec:lei}

A Lei Estadual nº 16.820 foi sancionada em 8 de janeiro de 2019 e entrou em
vigor em 9 de janeiro. Ela acrescenta o art. 28-B à Lei Estadual nº
12.228/1993, a norma estadual sobre agrotóxicos, com a seguinte redação:

\begin{citacao}
Art. 28-B. É vedada a pulverização aérea de agrotóxicos na agricultura no
Estado do Ceará.

§ 1º A infração ao art. 1º sujeita o infrator ao pagamento de multa de 15 mil
(quinze mil) UFIRCEs.

§ 2º Fica proibida a incorporação de mecanismos de controle vetorial por meio
de dispersão por aeronave em todo o Estado do Ceará, inclusive para os casos de
doenças causadas por vírus.
\end{citacao}

Dois traços da redação organizam o desenho empírico. A proibição é de método:
nenhum princípio ativo é nomeado. E ela não se restringe à agricultura, porque
o § 2º alcança a dispersão aérea para controle de vetores. Um município sem
agricultura pulverizada, mas com controle aéreo de endemias, também é alcançado
pela norma.

A Tabela~\ref{tab:cronologia} reúne a cronologia. Ela distingue três marcos de
antecipação: a notícia (apresentação do projeto de lei), a certeza (aprovação
unânime) e a obrigação (vigência). A distinção tem consequência empírica,
discutida na Seção~\ref{sec:dose}: a janela em que se mede a exposição anterior
à lei começa depois do primeiro marco.

\begin{table}[htb]
\centering
\caption{Cronologia legal do tratamento}
\label{tab:cronologia}
\small
\begin{tabular}{lp{11cm}}
\toprule
Data & Evento \\
\midrule
09/12/1993 & Lei Estadual nº 12.228/1993, norma estadual sobre agrotóxicos. \\
20/11/2009 & Limoeiro do Norte proíbe a pulverização aérea pela Lei Municipal
             nº 1.478/2009. \\
26/05/2010 & A Lei Municipal nº 1.511/2010, de política ambiental, revoga a
             proibição de Limoeiro do Norte, um mês após o assassinato de Zé
             Maria do Tomé. \\
24/02/2015 & Apresentação do Projeto de Lei nº 18/2015 (marco de notícia). \\
18/12/2018 & Aprovação por unanimidade na Assembleia Legislativa (marco de
             certeza). \\
09/01/2019 & Vigência da Lei nº 16.820/2019 (marco de obrigação). \\
mai/2023   & O STF julga a lei constitucional (ADI 6137). \\
19/12/2024 & A Lei Estadual nº 19.135/2024 abre exceção para aeronaves
             remotamente pilotadas. \\
2025       & Nova contestação no STF (ADI 7794), pendente. \\
\bottomrule
\end{tabular}
\legend{Fonte: elaboração própria, a partir dos textos oficiais das Leis
Estaduais nº 12.228/1993 e nº 16.820/2019 e do acervo da Câmara Municipal de
Limoeiro do Norte. A revogação de 2010 está documentada por fonte secundária
(imprensa); o texto integral da Lei nº 1.511/2010 ainda será conferido.}
\end{table}

A lei de 2024 mantém a vedação, salvo quando a aplicação é feita por aeronave
remotamente pilotada, sob condições de altura, velocidade do vento e distância
de escolas, hospitais e áreas protegidas. A nova redação não reproduz o § 2º de
2019. Do ponto de vista de instrumento, o Ceará não revogou a proibição:
substituiu a cota zero por regulação de desempenho com zona-tampão. Para o
Ensaio 1, a consequência é encerrar a janela de análise em 19 de dezembro de
2024, de modo que o tratamento avaliado seja a proibição tal como vigorou entre
2019 e 2024.

\section{A Chapada do Apodi e a química do caso}
\label{sec:chapada}

A Chapada do Apodi, no Baixo Jaguaribe, concentra a fruticultura irrigada de
exportação do Ceará, sobretudo banana e melão, em perímetros de irrigação sobre
solo cárstico. É a região a que a lei responde.

A justificativa do Projeto de Lei nº 18/2015 reproduz uma tabela de análises de
água atribuída ao Dossiê ABRASCO, referente a 23 pontos de coleta na Chapada em
2009 (Tabela~\ref{tab:abrasco}). O perfil é de fungicida e inseticida sobre
fruticultura, não de herbicida sobre grão transgênico.

\begin{table}[htb]
\centering
\caption{Princípios ativos detectados em água na Chapada do Apodi, 2009}
\label{tab:abrasco}
\small
\begin{tabular}{llc}
\toprule
Princípio ativo & Classe & Amostras positivas \\
\midrule
Procimidona  & Fungicida                    & 23 de 23 \\
Carbaril     & Inseticida (carbamato)       & 23 de 23 \\
Carbofurano  & Inseticida (carbamato)       & 18 de 23 \\
Fenitrotiona & Inseticida (organofosforado) & 16 de 23 \\
Glifosato    & Herbicida                    & 4 de 23 \\
\bottomrule
\end{tabular}
\legend{Fonte: tabela reproduzida na justificativa do Projeto de Lei nº
18/2015, atribuída ao Dossiê ABRASCO. A fonte primária do Dossiê ainda será
conferida.}
\end{table}

Disso decorrem três consequências para o projeto. Os desenhos construídos sobre
soja tolerante a herbicida transferem na forma, isto é, na ideia de usar
aptidão agroclimática como fonte de variação, e não no mecanismo. A lista de
culturas usada na literatura para construir o índice de aptidão, que é de
culturas transgênicas, precisa ser adaptada, porque banana não é transgênica. E
carbamatos e organofosforados são inibidores de colinesterase, com quadro agudo
característico, o que torna a notificação de intoxicação aguda um canal
informativo.

\section{O mecanismo: deriva e substituição}
\label{sec:deriva}

O que distingue a aplicação aérea da terrestre é a deriva. A exposição
resultante é difusa e alcança a população, não apenas o trabalhador. Dois
canais operam em paralelo: o atmosférico, por deposição direta sobre quem está
a favor do vento, e o hídrico, por escoamento e percolação até a água de
consumo.

Há um terceiro canal, de sinal oposto. Proibido o avião, a aplicação não
desaparece: migra para trator e pulverizador costal. Isso afasta o produto da
população e o aproxima do aplicador. A proibição pode reduzir a exposição
populacional e aumentar a intoxicação aguda ocupacional, e é por isso que o
projeto mede esse canal em separado (hipóteses H3 e H4).

\section{Fiscalização e registros administrativos}
\label{sec:fiscalizacao}

A fiscalização do art. 28-B cabe, pela Lei nº 12.228/1993, à Superintendência
Estadual do Meio Ambiente, à Secretaria da Agricultura e à Secretaria da Saúde.
A extensão efetiva da fiscalização não é conhecida. Se a proibição não foi
fiscalizada, a exposição anterior mede intenção de aplicação, e o parâmetro
estimado é um efeito de intenção de tratar.

Três registros administrativos podem dizer quem aplicava por aeronave antes da
lei, e nenhum deles foi obtido até aqui. O primeiro é o cadastro estadual de
empresas prestadoras de serviço (art. 8º da Lei nº 12.228/1993), a cargo da
SEMACE. O segundo é o receituário agronômico, que registra a modalidade de
aplicação (Decreto Federal nº 4.074/2002, art. 66), a cargo da ADAGRI. O
terceiro são os relatórios mensais que os operadores aeroagrícolas devem
entregar à Superintendência Federal de Agricultura do estado em que atuaram,
por força da Instrução Normativa MAPA nº 2/2008. As minutas dos três pedidos de
acesso à informação estão redigidas, e o protocolo é a primeira atividade do
plano de trabalho (Capítulo~\ref{cap:plano}).

% =====================================================================
% 3 Referencial teórico e revisão da literatura
% Fontes: paper/secoes/03-teoria.tex; paper/secoes/01-introducao.tex
% (antecedentes); docs/ars/13 (Calzada); docs/ars/16 (Denteh & Kédagni).
% =====================================================================

\chapter{Referencial teórico e revisão da literatura}
\label{cap:referencial}

\section{A deriva como externalidade negativa}
\label{sec:externalidade}

O agrotóxico que deriva sobre o vizinho é um custo imposto a terceiros que não
passa pelo sistema de preços. O produtor que contrata a aplicação aérea
internaliza o custo do serviço e do produto, e não o custo em saúde suportado
pela população a favor do vento. O resultado privado é aplicação acima do nível
socialmente eficiente.

As duas respostas canônicas divergem no instrumento. A pigouviana internaliza o
custo por tributo igual ao dano marginal. A coaseana observa que, com direitos
de propriedade bem definidos e custos de transação baixos, as partes negociam a
solução eficiente \cite{coase1960}, e desloca a pergunta para quem detém o
direito: o de pulverizar ou o de não ser pulverizado. A barganha coaseana não
se aplica a este caso. A deriva atinge população rural difusa, numerosa e não
organizada, com dano probabilístico, latente e de difícil atribuição
individual. Os custos de transação são proibitivos, e é isso que justifica a
intervenção regulatória e leva à discussão sobre o instrumento.

\section{Escolha de instrumento}
\label{sec:instrumento}

Sob incerteza a respeito dos custos de controle, um instrumento de preço e um
de quantidade não são equivalentes em bem-estar esperado \cite{weitzman1974}.
Na aproximação linear-quadrática, qual deles domina depende das inclinações
relativas do benefício marginal e do custo marginal do controle: se o benefício
marginal é mais inclinado, domina a quantidade; se o custo marginal é mais
inclinado, domina o preço.

Três qualificações da literatura importam para este caso, e o
Capítulo~\ref{cap:ensaio2} as desenvolve. A primeira é que proibir um método de
aplicação é um padrão tecnológico, e padrões sobre insumo ou tecnologia têm
efeitos distintos dos de padrões sobre emissão \cite{helfand1991}. A segunda é
que a deriva é poluição de fonte difusa: o regulador não observa a contribuição
de cada produtor, e o instrumento viável depende do que se consegue monitorar
\cite{segerson1988}. Sob heterogeneidade entre produtores, a taxa tende a ser
mais eficiente e a regulação por quantidade tende a gerar maior lucro ao setor
regulado \cite{helfand1995}. A terceira é que a literatura específica sobre
pesticidas trata dos custos de bem-estar e dos efeitos distributivos da
regulação \cite{lichtenberg1988,fernandezcornejo1998} e registra cenários em
que cotas superam taxas \cite{skevas2012}, além de simulação de banimento de
molécula na Europa \cite{bocker2020glyphosate}.

Há uma camada normativa que o cálculo de bem-estar esperado não capta.
Precificar uma externalidade de saúde supõe que o dano é comensurável com
dinheiro e que a compensação é resposta adequada. Proibi-la supõe um direito a
não ser exposto, anterior ao cálculo. \citet{hitzig2020normative} chama de
\emph{normative gap} a distância entre os objetivos de quem formula a política
e os objetivos que o economista embute no desenho do instrumento. A lei
cearense e a decisão do STF que a validou raciocinam em termos de saúde,
ambiente e exposição involuntária de terceiros. Ordenar proibição e taxa apenas
por bem-estar esperado responde a uma pergunta que o legislador não formulou.

\section{Evidência empírica sobre agrotóxicos e saúde ao nascer}
\label{sec:evidencia}

A evidência econômica se organiza em torno de duas perguntas: o que uma
substância faz à saúde e o que um método de aplicação faz à saúde.

Do lado da substância, \citet{dias2023down} são a referência brasileira.
Combinam o ganho potencial de produtividade trazido pela soja transgênica a
cada município com a direção do fluxo d'água na bacia e documentam piora dos
desfechos ao nascer a jusante. \citet{reynier2025glyphosate} estimam, para os
Estados Unidos, redução de 23 a 32 g no peso ao nascer à intensidade média de
exposição, concentrada no decil inferior de peso esperado.
\citet{larsen2017agricultural} encontram, no Vale de San Joaquin, efeitos
adversos de 5 a 9\% apenas entre os expostos ao topo da distribuição, o que
recomenda tratar a exposição como contínua e desconfiar de linearidade.
\citet{frank2024ecosystem} exploram um choque que eleva o uso de inseticida, a
perda de morcegos predadores, e encontram aumento da mortalidade infantil nos
condados atingidos.

Do lado do método, \citet{camacho2017health} avaliam as consequências
sanitárias da aspersão aérea de herbicida sobre cultivos ilícitos na Colômbia.
O análogo químico e de cultura mais próximo é \citet{calzada2023bananas}, sobre
a fumigação aérea de bananais no Equador, com déficit de peso ao nascer de 80 a
150 g sob exposição alta, segundo o resumo do artigo. Esse intervalo é a
referência de magnitude usada no exercício de calibração da
Seção~\ref{sec:diluicao}.

Nenhum desses trabalhos avalia uma proibição de método. É essa a lacuna que o
Ensaio 1 ocupa.

\section{Literatura metodológica}
\label{sec:metodo-lit}

A proibição entrou em vigor de uma só vez em todo o estado. Não há variação na
data de adoção, e os estimadores desenhados para adoção escalonada
\cite{callaway2021multiple,goodmanbacon2021,dechaisemartin2020} não se aplicam.
A variação disponível é de intensidade. \citet{callaway2024continuous} mostram
que diferenças em diferenças com tratamento contínuo não recupera a curva
dose-resposta sob as hipóteses habituais e distinguem dois parâmetros: o efeito
no nível de dose recebido, identificado sob tendências paralelas usuais, e a
curva causal, que exige tendências paralelas fortes.

Quatro outras peças sustentam o desenho. \citet{roth2023trending} mostra que um
teste de tendências anteriores pode passar por falta de poder.
\citet{rambachan2023credible} oferecem intervalos que admitem violações
limitadas do paralelismo. \citet{conley2011} e \citet{cameron2008wild} tratam
da inferência com poucos grupos, com as ressalvas de \citet{canay2021wild}. E
\citet{denteh2022misclassification} estudam diferenças em diferenças com
tratamento mal classificado: quando não há falso negativo, o parâmetro estimado
é o efeito verdadeiro multiplicado pelo valor preditivo positivo da
classificação. Esse resultado organiza a leitura dos resultados preliminares
(Seção~\ref{sec:diluicao}).

% =====================================================================
% 4 Metodologia do Ensaio 1
% Fontes: paper/secoes/04-identificacao.tex, 05-dados.tex;
% docs/pre-especificacao.md; docs/ars/17 (resposta ao parecer).
% Posição já corrigida pelo parecer de 22/09/2026 — ver pauta do orientador.
% =====================================================================

\chapter{Metodologia do Ensaio 1}
\label{cap:metodologia}

A exposição segue a ordem parâmetro-alvo, hipóteses de identificação e
estimação. As escolhas descritas neste capítulo são as do plano de análise
fechado em 21 de setembro de 2026, com os desvios posteriores registrados no
Apêndice~\ref{ap:preespecificacao}.

\section{Parâmetro-alvo}
\label{sec:estimando}

A fonte de variação é a intensidade de exposição anterior à lei entre os 184
municípios cearenses. O tratamento é contínuo, e o estimador é o de
\citet{callaway2024continuous}, que separa dois parâmetros:

\begin{alineas}
  \item $\ATT$, o efeito médio, entre os municípios com dose $d$, de ter
    recebido $d$ em vez de zero. É um efeito no nível da dose e é identificado
    sob tendências paralelas no resultado não tratado;
  \item $ATE(d)$ e sua derivada $ACR(d)$, a curva dose-resposta causal. A curva
    exige tendências paralelas fortes (SPT), que restringem trajetórias sob
    doses que o município não recebeu e excluem seleção por ganhos.
\end{alineas}

O plano de análise declarou a curva como alvo primário. A estimação mostrou que
a derivada agregada não é interpretável (Seção~\ref{sec:res-nivel}), e desde 22
de setembro de 2026 o resultado principal é o $\ATT$. O desvio está registrado
com data e motivo.

A variável de dose é área plantada de banana, e não aplicação aérea. O que o
desenho estima é, portanto, o \textbf{efeito de forma reduzida da proibição
segundo uma \emph{proxy} de intensidade da bananicultura}. Ele não é, sem
hipóteses adicionais, o efeito da retirada da pulverização aérea onde ela
ocorria. \decisao{1}{declarar o estimando como forma reduzida segundo a
\emph{proxy}; a alternativa, dose-resposta da pulverização aérea, exigiria o
cadastro operacional anterior à lei.}

Sobre a hipótese forte, o teste de tendências anteriores não a isola, porque a
implicação testável antes do tratamento é comum às duas versões do paralelismo.
Ele pode, porém, falsificá-la: \citet{callaway2024continuous} propõem verificar
se a relação entre a variação do desfecho e a dose é nula no período anterior,
e o projeto implementa essa verificação (Seção~\ref{sec:res-robustez}).

\section{Variável de tratamento}
\label{sec:dose}

A dose é a área plantada média de 2015 a 2018 da cultura-âncora, na Produção
Agrícola Municipal do IBGE, normalizada. Usa-se área plantada, e não colhida,
porque o período de 2012 a 2017 foi de seca severa no semiárido e a área
colhida responde à quebra de safra.

A cultura-âncora foi escolhida pelo dado, e não por suposição. O melão tem área
positiva em 10 municípios e o algodão em 28, suporte insuficiente para a
estimação não paramétrica. A banana tem área positiva em 169 municípios, com
índice de Gini de 0,86 e 78\% da área no decil superior. É também a cultura que
corresponde à química documentada: a procimidona é fungicida de bananal.

Três limitações de medida acompanham a escolha.

\begin{alineas}
  \item A proibição é de método. Área plantada mede intensidade agrícola, e
    culturas e municípios diferem em quanto recebem aplicação por aeronave. A
    Seção~\ref{sec:diluicao} mostra que essa distância é grande.
  \item A janela de medida começa depois da notícia da proibição. O projeto de
    lei foi apresentado em 24 de fevereiro de 2015, e a janela de 2015 a 2018 é
    posterior a esse marco. Se produtores responderam à expectativa, a dose já
    incorpora parte do ajuste. \decisao{3}{manter 2015--2018 como janela
    principal e reportar 2010--2014 como sensibilidade; promover 2010--2014 a
    principal muda o plano de análise e pode mudar o grupo tratado.}
  \item O zero da PAM é zero de \emph{proxy}. O § 2º do art. 28-B alcança o
    controle vetorial por aeronave, de modo que um município sem banana pode
    ter sido tratado (Seção~\ref{sec:comparacao}).
\end{alineas}

\section{Exposição gestacional}
\label{sec:exposicao}

Um desfecho perinatal não é contemporâneo à exposição. Um nascimento de março
de 2019 foi gestado, em dois terços, antes da vigência da lei. O painel calcula,
para cada célula município-mês, a fração da janela gestacional posterior à
proibição. Apenas as coortes nascidas a partir de outubro de 2019 têm gestação
integralmente posterior. A janela gestacional é fixa: usar a duração observada
condicionaria a exposição a um desfecho do estudo, a prematuridade. As coortes
cuja gestação atravessa a vigência são descartadas na estimação principal, o
que corresponde a 9,4\% das células. As séries de intoxicação aguda usam o mês
do evento, sem retroprojeção.

\section{Grupo de comparação}
\label{sec:comparacao}

O grupo de comparação é o dos municípios de dose nula, construído de quatro
maneiras: área nula da cultura-âncora; área nula de qualquer cultura candidata
a aplicação aérea; ausência de operador aeroagrícola registrado; e baixa
aptidão agroclimática. A definição primária é a quarta, e as demais entram como
sensibilidade. O grupo tem 15 municípios (14 na quarta definição).

Que esse grupo seja nunca tratado é hipótese, e não fato estabelecido. O Censo
Agropecuário de 2006 a apoia para a aplicação agrícola: nenhum município de
dose nula tinha estabelecimento com aplicação por aeronave. O Censo não mede o
controle vetorial aéreo, e não há fonte que o meça com suficiência (o SINAN
registra 89 notificações por agente de uso em saúde pública em oito anos). O
peso dessa hipótese é alto, porque os mesmos 15 municípios servem de controle a
todos os estimadores binários do projeto.

Duas verificações relativas ao grupo tratado foram concluídas. A varredura das
leis municipais anteriores a 2019 nos 17 municípios do decil superior encontrou
uma única proibição, a de Limoeiro do Norte, revogada em 2010; não há,
portanto, município já tratado na janela de 2015 a 2018. E o índice de aptidão
agroclimática FAO-GAEZ, construído para os 184 municípios pela receita de
\citet{reynier2025glyphosate}, é maior onde há mais banana (média de 0,37 no
decil superior contra 0,26 nos demais). \decisao{2}{tratar a aptidão como
construção do grupo de comparação e checagem de relevância; estimação por
variável instrumental fica como extensão declarada, porque a aptidão no
semiárido se correlaciona com disponibilidade de água, que afeta o desfecho
diretamente.}

\section{Dados}
\label{sec:dados}

A Tabela~\ref{tab:fontes} resume as fontes. O painel do Ensaio 1 é município
por mês, balanceado, com 184 municípios e 96 meses (2015 a 2022), num total de
17.664 células, das quais 17.654 têm nascimento registrado. Os desfechos são o
peso ao nascer, o baixo peso (menos de 2.500 g) e a prematuridade (menos de 37
semanas), agregados à célula. A Tabela~\ref{tab:descritivas} traz as
descritivas. \decisao{4}{estender a série do SINASC até 19/12/2024, que já
está disponível na fonte original, refazer painel e análise de poder e
congelar a data de extração de cada base.}

\begin{table}[htb]
\centering
\caption{Fontes de dados e estado de aquisição}
\label{tab:fontes}
\small
\begin{tabular}{p{3.6cm}p{2.1cm}p{8.0cm}}
\toprule
Fonte & Estado & Uso e observação \\
\midrule
PAM/IBGE         & adquirida   & Dose: área plantada por cultura, 2015--2018 \\
SINASC/DATASUS   & adquirida   & Desfechos de nascimento, 2015--2022 \\
SIM/DATASUS      & adquirida   & Óbito fetal \\
SINAN            & adquirida   & Intoxicação exógena; 1.217 notificações por
                                 agrotóxico de uso agrícola \\
População (IBGE) & adquirida   & Denominador; 2022 vem do Censo \\
FAO-GAEZ         & adquirida   & Aptidão agroclimática, 184 municípios \\
Censo Agropecuário 2006 & adquirida & Estabelecimentos por equipamento de
                                 aplicação, incluída a aeronave \\
Leis municipais  & levantadas  & 17 municípios do decil superior \\
SISAGUA          & suspensa    & Quebra de registro laboratorial no Ceará a
                                 partir de 2020 \\
ANA (bacias)     & suspensa    & Depende de coordenada que o SISAGUA não traz \\
INMET            & suspensa    & Direção do vento sem variação espacial \\
MapBiomas        & descartada  & Banana em classe genérica \\
SEMACE, ADAGRI, MAPA & a solicitar & Registros de aplicação aérea anteriores
                                 à lei, por pedido de acesso à informação \\
\bottomrule
\end{tabular}
\legend{Fonte: elaboração própria. Estado em setembro de 2026.}
\end{table}

\begin{table}[htb]
\centering
\caption{Descritivas do painel, Ceará, 2015--2022}
\label{tab:descritivas}
\small
\begin{tabular}{lr}
\toprule
Grandeza & Valor \\
\midrule
Nascidos vivos (após limpeza)                 & 1.001.709 \\
Peso médio ao nascer                          & 3.208 g \\
Baixo peso ($<$ 2.500 g)                      & 8,34\% \\
Prematuridade ($<$ 37 semanas)                & 12,58\% \\
Óbitos fetais                                 & 11.240 \\
Notificações por agrotóxico agrícola (SINAN)  & 1.217 \\
Municípios com dose positiva                  & 169 \\
Municípios com dose nula                      & 15 \\
\bottomrule
\end{tabular}
\legend{Fonte: elaboração própria, a partir de SINASC, SIM, SINAN e PAM.}
\end{table}

Três fontes foram obtidas, inspecionadas e suspensas, e o motivo de cada
suspensão é informação sobre o desenho. No SISAGUA, três das quatro moléculas
mais frequentes na Chapada não constam da lista monitorada, e as detecções de
agrotóxico no Ceará caem de 44 em 2019 para zero em 2020, 2021 e 2022, com mais
amostras coletadas; um modelo de diferenças em diferenças sobre essa variável
atribuiria à lei uma mudança de prática laboratorial. Sem o SISAGUA, o desenho
de montante e jusante fica sem insumo. No INMET, a coerência de direção do
vento entre estações é de 0,96: os alísios sopram de leste, e estar a favor do
vento equivale a estar a oeste da fonte. O canal de intoxicação, previsto com
dados de internação, passou ao SINAN, porque a série hospitalar registra um
único caso de intoxicação acidental por agrotóxico no período.

\section{Estimação e inferência}
\label{sec:estimacao}

A estimação principal usa o estimador de \citet{callaway2024continuous} com
aproximação não paramétrica, sobre o painel colapsado em um período anterior e
um posterior. A implementação impõe três restrições que afetam o que se pode
afirmar: opera sobre dois períodos; não produz o estudo de eventos e a curva
pelo mesmo procedimento; e não acomoda covariáveis.

O efeito agregado no nível é, na implementação usada, a diferença entre a
variação média do desfecho nos 169 municípios de dose positiva e a variação
média nos 15 de dose nula. Na inferência, portanto, o grupo escasso é o de
controle. Reportam-se, para o nível, o erro-padrão do \emph{bootstrap}
multiplicador, o \emph{wild bootstrap} com hipótese nula imposta
\cite{cameron2008wild}, uma análise de permutação e a exclusão sucessiva de
cada município de controle. A permutação da dose é reportada como análise de
permutação, e não como inferência exata: a exatidão exigiria que a dose tivesse
sido atribuída de forma intercambiável, o que não ocorre, porque a área de
banana segue aptidão e geografia. \decisao{5}{adotar como procedimento principal
para o nível uma inferência que respeite os 15 controles, na linha de
\citet{conley2011} com os papéis de tratado e controle invertidos.}

A camada de robustez compreende o poder do teste de tendências anteriores
\cite{roth2023trending}, a sensibilidade a violações de paralelismo
\cite{rambachan2023credible}, a verificação da inclinação em dose em datas
fictícias do período anterior e o controle sintético com diferenças em
diferenças \cite{arkhangelsky2021sdid}, na versão intraestadual. Uma
especificação com pareamento por escore de propensão e diferenças em diferenças
\cite{bravoureta2020psmdid} entra como camada de comunicação dos resultados.

\section{Poder, falsificação e multiplicidade}
\label{sec:poder}

O desenho declara seu poder antes de estimar. O efeito mínimo detectável
previsto para o peso ao nascer era de 33,4 g. O critério de falsificação fixa
um piso de efeito esperado de 15 g, abaixo do extremo inferior do intervalo de
23 a 32 g de \citet{reynier2025glyphosate}, por se supor atenuação quando se
remove um método e não uma molécula. Se a meia-largura do intervalo de
confiança ficar acima do piso, o resultado é reportado como limite informativo
sobre o parâmetro de forma reduzida, e não como achado.

A família confirmatória tem duas hipóteses, peso médio e óbito fetal, com
correção de Holm. Todo o restante é exploratório e declarado como tal. O plano
veda trocar o desfecho ou descer de unidade geográfica depois de ver o
resultado.

\section{Ameaças à identificação}
\label{sec:ameacas}

O Quadro~\ref{qua:ameacas} resume as ameaças e o tratamento dado a cada uma.

\begin{quadro}[htb]
\centering
\caption{Ameaças à identificação e resposta do desenho}
\label{qua:ameacas}
\small
\begin{tabular}{|p{3.6cm}|p{4.9cm}|p{5.6cm}|}
\hline
\textbf{Ameaça} & \textbf{Efeito sobre o resultado} & \textbf{Resposta} \\
\hline
Erro de classificação do tratamento
  & Atenua o efeito estimado; é o mecanismo da falta de poder
  & Medir a aplicação aérea por registro administrativo; concentrar a
    comparação onde havia aeronave \\
\hline
Prática interrompida antes de 2019
  & Se a pulverização cessou na Chapada em 2011, a lei não removeu exposição
  & Consultar a ação civil pública de 2011 e os relatórios do MAPA \\
\hline
Grupo de comparação exposto
  & Desloca o nível; pode inverter o sinal
  & Quatro construções do grupo; exclusão sucessiva dos controles \\
\hline
Antecipação desde 02/2015
  & Contamina a dose, não só o desfecho
  & Janela de dose 2010--2014 como sensibilidade; estudo de eventos até 2018 \\
\hline
Tendências paralelas fortes
  & A curva deixa de ser identificada
  & Resultado principal no nível; verificação em datas fictícias \\
\hline
Seleção para nascimento vivo
  & Efeito real pode aparecer nulo ou invertido
  & Óbito fetal como segundo desfecho confirmatório \\
\hline
Fiscalização desconhecida
  & O parâmetro é de intenção de tratar
  & Pedidos de acesso à informação; canal de intoxicação \\
\hline
Poucos grupos
  & Precisão aparente maior que a real
  & Inferência que respeita os 15 controles \\
\hline
\end{tabular}
\legend{Fonte: elaboração própria.}
\end{quadro}

% =====================================================================
% 5 Resultados preliminares e diagnóstico de poder
% Fontes: paper/secoes/06-diagnostico-poder.tex (números);
% docs/ars/17 §1 e §7 (leitura corrigida); docs/ars/16 (diluição);
% docs/ars/19 §2.3 (2011).
%
% Todo número deste capítulo vem da execução com dados reais de 21–22/09/2026
% e está transcrito de paper/secoes/06-diagnostico-poder.tex. Nenhum número
% novo foi produzido para este documento. Rastreio no Apêndice B.
% =====================================================================

\chapter{Resultados preliminares}
\label{cap:resultados}

A estimação do Ensaio 1 foi executada uma vez com dados reais, em setembro de
2026, sobre o painel de 2015 a 2022. Este capítulo reporta o que ela mostrou.
Os números são preliminares em dois sentidos: a série de nascimentos ainda será
estendida até 2024, e a inferência sobre o resultado principal ainda será
refeita com os procedimentos descritos na Seção~\ref{sec:estimacao}.

\section{Efeito no nível e curva dose-resposta}
\label{sec:res-nivel}

O $\ATT$ agregado para o peso ao nascer é de $-36{,}19$ g, com erro-padrão de
15,47 g pelo \emph{bootstrap} multiplicador e intervalo de 95\% em
$[-66{,}51;\ -5{,}87]$, que exclui o zero. Trocar a construção do grupo de
comparação desloca a estimativa para $-35{,}33$ g. O sinal é o contrário do
previsto pela hipótese H1: a estimativa aponta redução do peso ao nascer nos
municípios com banana em relação aos sem banana, depois da lei. A hipótese H2
admite esse sinal, mas o resultado do óbito fetal não a sustenta
(Seção~\ref{sec:res-inclinacao}).

Três razões impedem tratar esse número como achado, e nenhuma delas é
compatibilidade com o efeito nulo.

A primeira é o critério fixado antes da estimação. A meia-largura do intervalo
é de 30,32 g, o dobro do piso de 15 g.

A segunda é a inferência. O agregado é a diferença entre a variação média dos
169 municípios com área positiva de banana e a dos 15 de área nula. Com quinze
controles, o \emph{bootstrap} multiplicador é otimista. A inferência que
respeita esse número de controles está implementada e ainda não foi executada
com os dados reais; é a primeira tarefa computacional do plano de trabalho.

A terceira é uma calibração, apresentada na Seção~\ref{sec:diluicao}: lido como
efeito diluído da retirada da pulverização aérea, o coeficiente implicaria uma
magnitude implausível entre os municípios de fato tratados.

A curva dose-resposta não fornece resultado utilizável. A derivada estimada é
fortemente não linear e muda de sinal ao longo da dose. Na faixa de dose
inferior a 0,1\%, onde estão 48 dos 169 pontos, ela é da ordem de $-9.100$ g
por unidade de dose, com erro-padrão maior que a estimativa. O agregado que o
estimador reporta é média simples sobre os pontos e é dominado por essa região.
Nenhum ponto exclui o zero. Isso é imprecisão sob hipóteses mantidas, e não
demonstração de que a inclinação seja irrestrita.

\section{Inclinação em dose e família confirmatória}
\label{sec:res-inclinacao}

A inclinação do desfecho em relação à dose é outro estimando, distinto do
nível. A estimativa é de $+6{,}81$ g por unidade de dose normalizada, e nenhum
dos três procedimentos de inferência a distingue de zero
(Tabela~\ref{tab:inferencia}).

\begin{table}[htb]
\centering
\caption{Inferência sobre a inclinação em dose, peso ao nascer}
\label{tab:inferencia}
\small
\begin{tabular}{lr}
\toprule
Procedimento & Valor-$p$ \\
\midrule
Erro-padrão agrupado por município     & 0,736 \\
\emph{Wild bootstrap} por município    & 0,757 \\
Análise de permutação da dose          & 0,793 \\
\bottomrule
\end{tabular}
\legend{Fonte: elaboração própria. O intervalo por inversão do teste de
permutação será recalculado, porque o valor obtido em setembro usava uma
aproximação.}
\end{table}

Um deslocamento de nível entre municípios com e sem banana, sem gradiente entre
as doses positivas, é mais compatível com uma diferença de trajetória entre os
dois grupos do que com uma resposta crescente à intensidade. O perfil da
variação do desfecho por faixa de dose, já implementado, é o diagnóstico
descritivo que mostrará onde está a variação.

Na família confirmatória, com correção de Holm, os valores-$p$ do peso médio
passam de 0,394 para 0,427 no teste unilateral e de 0,794 para 0,866 no
bilateral; os do óbito fetal, de 0,214 para 0,427 e de 0,433 para 0,866.
Nenhuma hipótese é rejeitada a 5\%. O coeficiente do óbito fetal é próximo de
zero, e 98,7\% das células têm menos de cinco eventos, de modo que a hipótese
H2 não é sustentada nem excluída.

\section{Robustez}
\label{sec:res-robustez}

\subsection{Poder do teste de tendências anteriores}

O teste conjunto dos coeficientes anteriores à lei tem poder de 80\% contra uma
tendência linear de 0,238 g por mês, que acumularia 11,4 g em quatro anos de
período posterior, e de 90\% contra 0,429 g por mês (20,6 g). O teste não
rejeitou. Isso é compatível com tendências menores que essas e, em uma de cada
cinco amostras, com tendências desse tamanho. Nada diz sobre violações não
lineares. Onze gramas é da ordem do efeito procurado.

\subsection{Sensibilidade a violações de paralelismo}

A análise de \citet{rambachan2023credible} foi aplicada a um estudo de eventos
com tratamento binário, que compara o decil superior de dose com o grupo de
dose nula. Sob paralelismo exato, o intervalo é $[-72{,}5;\ +24{,}5]$ g.
Admitida violação do tamanho da maior violação observada antes da lei, ele
passa a $[-217{,}4;\ +161{,}4]$ g. O resultado condiciona à hipótese de que o
grupo de dose nula não foi tratado.

\subsection{Inclinação em dose em datas fictícias}

Simulou-se a proibição em três datas do período anterior e estimou-se a mesma
inclinação em dose (Tabela~\ref{tab:spt}). Nenhuma das três rejeita. Para o
peso, porém, as três inclinações fictícias são maiores em módulo que a do
desenho real. O ruído do período anterior é maior que o efeito estimado: o
desenho não separa efeito de ruído para o peso ao nascer. Para o óbito fetal, a
relação se inverte, e nenhuma data fictícia alcança a inclinação real.

\begin{table}[htb]
\centering
\caption{Inclinação em dose em datas fictícias do período anterior}
\label{tab:spt}
\small
\begin{tabular}{lrrrr}
\toprule
& \multicolumn{2}{c}{Peso médio} & \multicolumn{2}{c}{Óbito fetal} \\
\cmidrule(lr){2-3}\cmidrule(lr){4-5}
Corte & Inclinação & Valor-$p$ & Inclinação & Valor-$p$ \\
\midrule
jan/2016 a nov/2016 & 25,76 & 0,371 & 0,00086 & 0,864 \\
jul/2016 a mai/2017 & 31,07 & 0,239 & 0,00074 & 0,852 \\
jan/2017 a nov/2017 & 12,96 & 0,608 & 0,00021 & 0,959 \\
\midrule
Desenho real        & 6,81  & 0,794 & $-0{,}00244$ & 0,433 \\
\bottomrule
\end{tabular}
\legend{Fonte: elaboração própria.}
\end{table}

\subsection{Controle sintético}

Três estimadores foram aplicados à mesma matriz, comparando municípios de dose
alta com os de dose nula: controle sintético com diferenças em diferenças
($-20{,}16$ g; intervalo $[-56{,}43;\ +16{,}10]$), controle sintético
($-34{,}64$ g; $[-67{,}43;\ -1{,}85]$) e diferenças em diferenças
($-21{,}85$ g; $[-59{,}54;\ +15{,}84]$). Um dos três exclui o zero. Ele não é
tratado como achado: os três usam os mesmos dados, nenhum é confirmatório e
escolher o que cruzou a linha é a prática que o plano de análise veda.

Todos os estimadores binários deste capítulo usam os mesmos 15 municípios de
controle. A concordância de sinal entre eles não equivale a evidências
independentes.

\section{O que a variável de dose mede: erro de classificação do tratamento}
\label{sec:diluicao}

O Censo Agropecuário de 2006 registra o número de estabelecimentos que usavam
agrotóxicos por tipo de equipamento, com a aeronave como categoria própria. No
Ceará eram 36 estabelecimentos, 0,03\% dos que usavam agrotóxicos, distribuídos
por sete municípios. Limoeiro do Norte (18) e Quixeré (9), ambos na Chapada do
Apodi, respondiam por 75\%.

Confrontada com esse dado, a dose por área plantada classifica mal o
tratamento. Dos 17 municípios do decil superior de banana, 15 não tinham
nenhum estabelecimento com aplicação por aeronave. Dos 169 municípios com dose
positiva, 7 tinham. Nenhum município de dose nula tinha.

Quando o erro é só de falso positivo, o resultado de
\citet{denteh2022misclassification} implica que o parâmetro estimado, $\theta$,
é o efeito verdadeiro multiplicado pelo valor preditivo positivo da
classificação: $\theta = VPP \cdot ATT$. A atenuação é pura, e o sinal não se
inverte. Com o VPP de 2 em 17 no decil superior e o efeito do análogo mais
próximo, de 80 a 150 g \cite{calzada2023bananas}, o efeito esperado no
parâmetro estimado é de no máximo cerca de 9 g na calibração do projeto (18 g
no caso-limite em que toda a população desses municípios estivesse exposta ao
efeito máximo), e o poder do desenho contra 9 g é de no máximo 11\%. A diluição é, assim, o mecanismo da falta de poder. Aumentar o número de
municípios tratados baixando o corte de dose pioraria o problema, porque os
municípios seguintes na ordenação não tinham aviação agrícola.

A mesma identidade serve para calibrar a leitura do $-36{,}19$ g. Suponha-se
que o Censo de 2006 classifique corretamente quem pulverizava por avião na
véspera da lei e que não haja exposição omitida no grupo de comparação. Sob
essas hipóteses fortes, o VPP no agregado é de 7 em 169, e o coeficiente
implicaria um efeito de cerca de $-874$ g entre os municípios de fato tratados,
seis vezes o maior efeito do análogo. O cálculo não identifica o efeito
verdadeiro, porque o Censo de 2006 não estabelece quem era tratado em 2018. Ele
mostra que o coeficiente dificilmente é uma versão atenuada de um efeito
sanitário plausível. Restam quatro explicações, que os dados atuais não
separam: o Censo deixou de representar a distribuição da aviação agrícola; há
diferença de trajetória entre municípios com e sem banana; a inferência
subestima a incerteza vinda dos 15 controles; ou houve exposição omitida no
grupo de comparação, por controle vetorial aéreo, caso em que a identidade não
vale e o sinal pode se inverter.

Há um risco adicional, aberto em setembro de 2026. Em março de 2011, noticiou-se
que as empresas de fruticultura da Chapada do Apodi estavam proibidas de
pulverizar por aeronave, e em julho do mesmo ano o Ministério Público Federal,
o do Trabalho e o estadual ajuizaram ação civil pública pedindo a proibição. O
desfecho de ambos é desconhecido. Se a prática cessou na Chapada em 2011, a lei
de 2019 não removeu exposição nos dois municípios que concentravam a aviação
agrícola, e o resultado refletiria ausência de tratamento, e não falta de
poder.

\section{Síntese}
\label{sec:sintese}

O projeto especificou um desenho de avaliação causal, fechou o plano de análise
antes de estimar e executou a estimação e a camada de robustez. O que os
resultados preliminares estabelecem é o seguinte.

\begin{alineas}
  \item O efeito mínimo detectável previsto (33,4 g) e a meia-largura realizada
    (30,3 g) coincidem. Isso valida o diagnóstico amostral do desenho. Não
    valida a medida do tratamento.
  \item O parâmetro estimado é o efeito de forma reduzida da proibição segundo
    a área de banana. O intervalo obtido limita esse parâmetro e não limita o
    efeito onde havia pulverização aérea.
  \item A estimativa no nível é negativa e seu intervalo assintótico exclui o
    zero, mas ela não satisfaz o critério de precisão, depende de 15 controles
    e não é plausível como efeito diluído da lei.
  \item A curva dose-resposta não é informativa com a precisão obtida.
  \item A restrição ativa é a medida do tratamento. O próximo passo é obtê-la e
    testar um desenho que concentre a comparação onde havia aeronave.
\end{alineas}

% =====================================================================
% 6 Ensaio 2 — proposta
% Fontes: docs/ensaio2-rota-inversao.md; docs/ars/17 §2 (comentários 1, 2, 6),
% §6 (Hitzig) e §7.6; docs/legislacao/adi-5553-7755-tributacao.md.
% Este capítulo é PLANO, não resultado. O enquadramento é a Decisão D6.
% =====================================================================

\chapter{Ensaio 2: escolha de instrumento}
\label{cap:ensaio2}

\section{Pergunta e ponto de partida}
\label{sec:e2-pergunta}

O Ensaio 2 pergunta se a proibição é o instrumento adequado para a
externalidade da pulverização aérea, em comparação com a taxação do insumo ou
do serviço e com a regulação de desempenho por zonas-tampão. O caso oferece os
três instrumentos em disputa real: a proibição de 2019, a regulação de
desempenho para drones adotada em 2024 e, no plano federal, a desoneração
tributária de agrotóxicos em julgamento no STF (ADI 5553 e ADI 7755), que é a
discussão do instrumento de preço.

A concepção original do ensaio era ancorar a regra de \citet{weitzman1974} na
curva dose-resposta estimada no Ensaio 1. Essa ancoragem não se sustenta, por
quatro razões que o projeto registra.

\begin{alineas}
  \item O $\ATT$ não é o benefício social do controle, e sua derivada não é a
    inclinação do benefício marginal.
  \item A dose do Ensaio 1 mede intensidade de bananicultura e não equivale à
    quantidade que o regulador escolhe, que seria área tratada por via aérea ou
    a exposição resultante.
  \item A proibição é uma cota fixada em zero, isto é, um padrão tecnológico
    \cite{helfand1991}. Ela só coincide com o instrumento de quantidade da
    regra de Weitzman se o ótimo esperado for a solução de canto, e a
    aproximação quadrática local é mais frágil justamente nesse ponto.
  \item O Ensaio 1 não produz, com a precisão obtida, informação utilizável
    sobre a inclinação do benefício marginal.
\end{alineas}

\section{Enquadramento proposto}
\label{sec:e2-enquadramento}

\decisao{6}{enquadrar o Ensaio 2 por suficiência informacional, padrão
tecnológico e objetivo do formulador, em lugar de ancorá-lo no $\ATT$; a
alternativa é adiar o ensaio até haver curvatura identificada.} O ensaio se
organiza em três partes.

\subsection{O que a regra exigiria}

A regra de Weitzman entra como referência para organizar a informação
necessária, e não como fórmula que receba os parâmetros do Ensaio 1. Aplicá-la
pede três objetos definidos sobre a mesma quantidade: a quantidade regulada, a
inclinação do benefício marginal social em relação a ela e a inclinação do
custo marginal agregado de reduzi-la. O ensaio especifica cada objeto para o
caso cearense, indica que dado o mediria e mostra quais deles a evidência
disponível informa. Um exercício de inversão, já implementado, calcula que
combinações de inclinações fariam o ordenamento entre preço e quantidade mudar;
sem calibração, ele não conclui por nenhum instrumento, e esse é o resultado a
reportar.

\subsection{O lado do custo}

A série histórica de custos de produção da banana da Conab traz a operação com
avião como item próprio. Nos sistemas que a utilizam, ela custa de R\$ 280 a
R\$ 520 por hectare e representa cerca de 1,5\% do custo total de produção.
Dos 88 sistemas acompanhados, 4 usam avião e 9 usam trator, e nenhum usa os
dois. Esses dados dão a ordem de
grandeza do custo da proibição. Não dão a curvatura do custo, que depende de
custos marginais dentro do produtor, de custos fixos, de capacidade e escala e
da heterogeneidade entre produtores, nem o custo contrafactual de converter
quem usava avião, porque a escolha de tecnologia é endógena ao sistema de
produção. O ensaio propõe levantar a dispersão do custo de conversão junto ao
setor e à literatura agronômica.

\subsection{O critério de comparação}

A lei declara um objetivo, que é exposição aérea nula. O ensaio propõe comparar
os instrumentos pelo custo e pela confiabilidade com que realizam esse
objetivo sob incerteza, em vez de ordená-los por bem-estar esperado. Sob esse
critério, uma taxa não garante exposição nula quando os custos são incertos, a
zona-tampão depende de fiscalização de desempenho, e a proibição tende a ter custo de
monitoramento menor, por ser observável a presença da aeronave. Essas são
proposições a examinar, e não conclusões. A literatura de
poluição difusa \cite{segerson1988,helfand1995} e a discussão sobre a distância
entre o objetivo do formulador e o do desenho econômico
\cite{hitzig2020normative} fornecem o referencial. A troca de instrumento de
2024, da cota zero para a regulação de desempenho com drones, é o caso a
analisar.

\section{Método e produto esperado}
\label{sec:e2-metodo}

O Ensaio 2 é analítico e documental. Usa o modelo de escolha de instrumento
sob incerteza como estrutura, os dados de custo da Conab como calibração de
ordem de grandeza, os resultados do Ensaio 1 como delimitação do que se sabe
sobre o benefício e os textos legais e as decisões do STF como fonte sobre os
objetivos do formulador. Os fundamentos do acórdão da ADI 6137 ainda serão
lidos na fonte.

O produto esperado é um ensaio que estabeleça (i) que informação a comparação
de instrumentos exige, (ii) que parte dela existe para o caso cearense e (iii)
como os três instrumentos se ordenam sob o critério de custo-efetividade e
confiabilidade, com as hipóteses de cada ordenamento explicitadas. O ensaio não
afirmará qual instrumento maximiza bem-estar esperado, porque a evidência não o
permite.

% =====================================================================
% 7 Plano de trabalho, riscos e contingências
% Fontes: docs/ars/19 §5 (matriz R1–R8) e §7; docs/ars/18;
% docs/legislacao/lai-*-minuta.md; docs/ars/17 §5.
% =====================================================================

\chapter{Plano de trabalho e riscos}
\label{cap:plano}

\section{Atividades}
\label{sec:atividades}

O trabalho restante se divide em cinco frentes.

\subsection{Medida do tratamento}

A restrição ativa do Ensaio 1 é saber onde e quando havia pulverização aérea
antes da lei. Três atividades a enfrentam.

\begin{alineas}
  \item Protocolar os pedidos de acesso à informação já minutados: ao MAPA,
    pelos relatórios mensais dos operadores aeroagrícolas com atuação no Ceará
    (o texto da Instrução Normativa nº 2/2008 será conferido antes do
    protocolo); à ADAGRI, pelo receituário agronômico com modalidade de
    aplicação; e à SEMACE, pelo cadastro de empresas prestadoras de serviço.
  \item Consultar, na Justiça Federal no Ceará, a ação civil pública de 2011
    contra empresas de fruticultura da Chapada do Apodi, para saber se houve
    liminar ou sentença que tenha interrompido a pulverização aérea antes de
    2019.
  \item Conferir o texto integral da Lei Municipal nº 1.511/2010 de Limoeiro do
    Norte.
\end{alineas}

\subsection{Consolidação dos dados e da inferência}

\begin{alineas}
  \item Estender o painel até 19 de dezembro de 2024 e registrar a data de
    extração de cada base;
  \item executar com dados reais a inferência sobre o nível que respeita os 15
    controles, o perfil por faixa de dose e o intervalo por inversão do teste
    de permutação, já implementados;
  \item reestimar a dose com a janela de 2010 a 2014 e medir a correlação de
    postos entre as duas janelas.
\end{alineas}

\subsection{Desenho com mais poder}

O projeto avaliou oito rotas de continuação. A mais promissora com os dados
existentes é um desenho de calendário: comparar, nos municípios onde havia
aeronave, a diferença de desfecho entre gestações expostas aos meses de pico de
aplicação e gestações expostas à entressafra, antes e depois da lei, contra a
mesma diferença nos demais municípios. A hipótese de identificação passa a ser
a estabilidade desse hiato sazonal, e não o paralelismo entre municípios com e
sem banana. Se os relatórios do MAPA forem obtidos, os municípios tratados e os
meses de pico deixam de ser supostos e passam a ser medidos.

\decisao{7}{adotar o desenho de calendário como etapa seguinte do Ensaio 1.
Ele muda a estratégia de identificação e exige plano de análise próprio,
fechado antes de cruzar exposição e desfecho.} A sequência prevista é: calcular
o efeito mínimo detectável do desenho de calendário apenas com dados anteriores
à lei, com o código já escrito; decidir com o orientador se o poder justifica a
adoção; fechar e versionar o plano de análise; e só então estimar.

\subsection{Ensaio 2}

Leitura do acórdão da ADI 6137 e dos autos das ADI 5553 e 7755; levantamento da
dispersão do custo de conversão de aplicação aérea para terrestre; redação das
três partes descritas na Seção~\ref{sec:e2-enquadramento}.

\subsection{Redação e defesa}

Revisão do texto do Ensaio 1 conforme o parecer recebido em setembro de 2026;
redação do Ensaio 2; introdução e conclusão gerais; depósito e defesa.

\section{Riscos e contingências}
\label{sec:riscos}

O Quadro~\ref{qua:riscos} relaciona os riscos principais.

\begin{quadro}[htb]
\centering
\caption{Riscos e contingências}
\label{qua:riscos}
\small
\begin{tabular}{|p{4.4cm}|p{4.2cm}|p{5.5cm}|}
\hline
\textbf{Risco} & \textbf{Consequência} & \textbf{Contingência} \\
\hline
Os pedidos de acesso à informação são negados ou os registros não existem
  & O tratamento continua medido por \emph{proxy}
  & Desenho de calendário com os municípios do Censo de 2006 e meses de pico
    declarados como hipótese \\
\hline
A pulverização aérea cessou na Chapada em 2011
  & Não há tratamento a avaliar onde ele se concentrava
  & O Ensaio 1 é apresentado como diagnóstico do desenho e da medida, com o
    achado documental como resultado \\
\hline
O desenho de calendário não tem poder
  & O Ensaio 1 permanece inconclusivo quanto ao efeito
  & Reportar o efeito de forma reduzida com inferência adequada e o limite de
    poder como contribuição metodológica \\
\hline
A série até 2024 altera as estimativas
  & Números do Capítulo~\ref{cap:resultados} mudam
  & As duas versões são reportadas, com data de extração \\
\hline
Prazo
  & Atraso do Ensaio 2
  & O Ensaio 2 é analítico e não depende de dado novo; pode ser redigido em
    paralelo \\
\hline
\end{tabular}
\legend{Fonte: elaboração própria.}
\end{quadro}

Em qualquer cenário, o Ensaio 1 tem um resultado defensável: o diagnóstico de
que o desenho disponível mede um efeito de forma reduzida, com o poder e a
qualidade da medida do tratamento quantificados contra critério fixado
previamente. Os cenários favoráveis acrescentam a esse diagnóstico uma
estimativa com tratamento medido.

\section{Reprodutibilidade e ética}
\label{sec:reprodutibilidade}

O código do projeto é versionado, com execução ordenada por um arquivo de
orquestração, ambientes fixados para Python e R e suíte de testes automatizada.
A preparação dos dados é feita em Python e a estimação com tratamento contínuo
em R, com troca de arquivos entre as duas etapas. O plano de análise e seus
desvios são versionados com data.

Os dados são secundários, públicos e anonimizados na origem (DATASUS, IBGE,
FAO). Os microdados não são redistribuídos, e os resultados são reportados em
agregados municipais. \pendente{confirmar com o programa se o projeto, por
usar apenas dados públicos sem identificação, dispensa apreciação por comitê
de ética, nos termos da Resolução CNS nº 510/2016.}

% =====================================================================
% 8 Cronograma
% ⚠️ O horizonte (dez/2026 a jul/2027) é SUPOSIÇÃO de trabalho, não dado do
% programa. Ajustar as colunas à data de defesa informada pelo pesquisador.
% =====================================================================

\chapter{Cronograma}
\label{cap:cronograma}

O Quadro~\ref{qua:cronograma} distribui as atividades do
Capítulo~\ref{cap:plano} por mês. \pendente{data prevista de defesa. O horizonte
abaixo, de dezembro de 2026 a julho de 2027, é suposição de trabalho e será
ajustado ao prazo do programa.} Os pedidos de acesso à informação têm prazo
legal de resposta de vinte dias, prorrogável por dez, e por isso são a primeira
atividade: o que eles devolverem condiciona o desenho de calendário.

\begin{quadro}[htb]
\centering
\caption{Cronograma de atividades}
\label{qua:cronograma}
\footnotesize
\setlength{\tabcolsep}{4pt}
\begin{tabular}{|p{7.4cm}|c|c|c|c|c|c|c|c|}
\hline
\textbf{Atividade} & \textbf{Dez} & \textbf{Jan} & \textbf{Fev} & \textbf{Mar}
  & \textbf{Abr} & \textbf{Mai} & \textbf{Jun} & \textbf{Jul} \\
 & 2026 & 2027 & 2027 & 2027 & 2027 & 2027 & 2027 & 2027 \\
\hline
Protocolo dos pedidos de acesso à informação (MAPA, ADAGRI, SEMACE)
  & X &   &   &   &   &   &   &   \\ \hline
Consulta à ação civil pública de 2011 e à Lei Municipal nº 1.511/2010
  & X &   &   &   &   &   &   &   \\ \hline
Extensão do painel até 2024 e congelamento das datas de extração
  & X & X &   &   &   &   &   &   \\ \hline
Inferência sobre o nível com 15 controles; perfil por dose; janela 2010--2014
  &   & X &   &   &   &   &   &   \\ \hline
Efeito mínimo detectável do desenho de calendário (dados anteriores à lei)
  &   & X &   &   &   &   &   &   \\ \hline
Tratamento das respostas aos pedidos de informação
  &   & X & X &   &   &   &   &   \\ \hline
Plano de análise do desenho de calendário, fechado e versionado
  &   &   & X &   &   &   &   &   \\ \hline
Estimação do desenho de calendário e robustez
  &   &   &   & X & X &   &   &   \\ \hline
Ensaio 2: leitura dos acórdãos e levantamento de custos
  &   & X & X & X &   &   &   &   \\ \hline
Redação do Ensaio 1 (versão revisada)
  &   &   &   & X & X & X &   &   \\ \hline
Redação do Ensaio 2
  &   &   &   &   & X & X &   &   \\ \hline
Introdução e conclusão gerais; revisão com o orientador
  &   &   &   &   &   & X & X &   \\ \hline
Depósito e defesa
  &   &   &   &   &   &   & X & X \\ \hline
\end{tabular}
\legend{Fonte: elaboração própria.}
\end{quadro}


% ----------------------------------------------------------------------
% Elementos pós-textuais
% ----------------------------------------------------------------------
\postextual

\bibliography{../paper/referencias}

\begin{apendicesenv}
\partapendices
% =====================================================================
% Apêndice A — resumo do plano de análise e dos desvios
% Fonte: docs/pre-especificacao.md (fechada em 2026-09-21, commit c386514).
% =====================================================================

\chapter{Plano de análise do Ensaio 1 e desvios registrados}
\label{ap:preespecificacao}

O plano de análise foi fechado e versionado em 21 de setembro de 2026. Ele não
é um pré-registro em sentido estrito: foi escrito depois dos diagnósticos de
desenho (distribuição da dose, dispersão das tendências municipais, efeito
mínimo detectável, relação entre aptidão e dose) e antes de qualquer estimativa
de efeito. A fronteira que ele protege é entre desenho e efeito. O
Quadro~\ref{qua:plano} resume as escolhas, e o Quadro~\ref{qua:desvios}, os
desvios posteriores.

\begin{quadro}[htb]
\centering
\caption{Escolhas do plano de análise}
\label{qua:plano}
\small
\begin{tabular}{|p{4.2cm}|p{10cm}|}
\hline
\textbf{Item} & \textbf{Escolha} \\
\hline
Desfecho primário & Peso médio ao nascer, por município e mês \\
\hline
Exposição primária & Fração da gestação posterior à vigência da lei \\
\hline
Cultura-âncora & Banana, por ter área positiva em 169 municípios e índice de
Gini da dose de 0,86 \\
\hline
Janela da dose & 2015--2018; 2010--2014 como sensibilidade \\
\hline
Grupo de comparação & Municípios de dose nula, definidos por baixa aptidão
agroclimática; área nula como sensibilidade \\
\hline
Cortes & Período anterior até dez/2018 e posterior a partir de out/2019;
coortes de transição descartadas (9,4\% das células) \\
\hline
Critério de falsificação & Piso de efeito esperado de 15 g \\
\hline
Família confirmatória & Peso médio e óbito fetal, com correção de Holm \\
\hline
Exploratório declarado & Exposição por trimestre; baixo peso e prematuridade;
canal de intoxicação; janela 2010--2014; definições alternativas do grupo de
comparação \\
\hline
Falsificação & Notificações de intoxicação intencional, que não devem se mover \\
\hline
Vedações & Não trocar o desfecho primário; não descer de unidade geográfica;
não reportar como achado coeficiente abaixo do efeito mínimo detectável; não
usar o canal hídrico nem o teste de direção do vento \\
\hline
\end{tabular}
\legend{Fonte: elaboração própria, a partir do plano de análise versionado.}
\end{quadro}

\begin{quadro}[htb]
\centering
\caption{Desvios do plano de análise}
\label{qua:desvios}
\small
\begin{tabular}{|p{2.1cm}|p{5.2cm}|p{6.9cm}|}
\hline
\textbf{Data} & \textbf{O que mudou} & \textbf{Motivo} \\
\hline
21/09/2026 & Grupo de comparação primário passa a ser o de baixa aptidão
  & O índice FAO-GAEZ foi obtido para os 184 municípios \\
\hline
22/09/2026 & Parâmetro principal passa da curva, $ACR(d)$, ao nível, $\ATT$
  & A derivada agregada é dominada pela faixa de dose em que o erro-padrão
    supera a estimativa \\
\hline
22/09/2026 & A hipótese de tendências paralelas fortes deixa de ser descrita
  como não testável
  & A verificação anterior ao tratamento não a isola, mas pode falsificá-la;
    foi implementada \\
\hline
22/09/2026 & A correção de Holm passa a ser calculada
  & Estava declarada e não implementada \\
\hline
\end{tabular}
\legend{Fonte: elaboração própria, a partir do plano de análise versionado.}
\end{quadro}

As decisões marcadas ao longo do texto como dependentes do orientador, se
adotadas, entrarão neste quadro com data e motivo, e os resultados originais
serão preservados.

% =====================================================================
% Apêndice B — de onde vem cada número
% =====================================================================

\chapter{Proveniência dos resultados preliminares}
\label{ap:proveniencia}

Os resultados do Capítulo~\ref{cap:resultados} vêm de uma única execução do
pipeline com dados reais, em 21 e 22 de setembro de 2026, na máquina do
pesquisador. Os microdados e os arquivos intermediários não são versionados. Os
números foram transcritos para o texto do Ensaio 1 e, dele, para este projeto;
nenhum número novo foi produzido para este documento. O
Quadro~\ref{qua:proveniencia} indica o programa que gera cada resultado, para
que a nova execução prevista no plano de trabalho possa ser conferida contra
esta.

\begin{quadro}[htb]
\centering
\caption{Programa de origem de cada resultado}
\label{qua:proveniencia}
\small
\begin{tabular}{|p{6.4cm}|p{7.8cm}|}
\hline
\textbf{Resultado} & \textbf{Programa (pasta \texttt{scripts/})} \\
\hline
Nível e curva dose-resposta & \texttt{estimate/03\_contdid.R} \\
\hline
Inclinação em dose e três procedimentos de inferência &
  \texttt{estimate/04\_robustness.py} \\
\hline
Poder do teste de tendências anteriores & \texttt{estimate/06\_pretrends.R} \\
\hline
Sensibilidade a violações de paralelismo & \texttt{estimate/07\_honestdid.R} \\
\hline
Controle sintético & \texttt{estimate/08\_synthdid.R} \\
\hline
Correção de Holm & \texttt{estimate/09\_holm.py} \\
\hline
Inclinação em datas fictícias & \texttt{estimate/10\_spt\_pretrend.py} \\
\hline
Erro de classificação e calibração &
  \texttt{estimate/12\_erro\_de\_classificacao.py} \\
\hline
Custos da Conab & \texttt{data\_prep/12\_clean\_conab\_custos.py} \\
\hline
Efeito mínimo detectável do desenho de calendário (a executar) &
  \texttt{estimate/14\_mde\_calendario.py} \\
\hline
\end{tabular}
\legend{Fonte: elaboração própria.}
\end{quadro}

\end{apendicesenv}

\end{document}
